\documentclass[11pt]{article}

% ---------------------------------------------------------------------------
% Alpha-GPT: a ground-up, optimization-first formulation (internal note)
% ---------------------------------------------------------------------------
\usepackage[margin=1in]{geometry}
\usepackage{amsmath,amssymb,amsthm}
\usepackage{booktabs}
\usepackage{enumitem}
\usepackage{microtype}
\usepackage{algorithm}
\usepackage{algpseudocode}
\usepackage{xcolor}
\usepackage[colorlinks=true,linkcolor=blue!50!black,citecolor=blue!50!black,urlcolor=blue!50!black]{hyperref}

\theoremstyle{definition}
\newtheorem{definition}{Definition}[section]
\theoremstyle{plain}
\newtheorem{proposition}{Proposition}[section]
\theoremstyle{remark}
\newtheorem{remark}{Remark}[section]

\newcommand{\E}{\mathbb{E}}
\newcommand{\R}{\mathbb{R}}
\newcommand{\Uni}{\mathcal{U}}
\newcommand{\Expr}{\mathcal{E}}
\newcommand{\Term}{\mathcal{T}}
\newcommand{\Op}{\mathcal{O}}
\newcommand{\Sharpe}{\mathrm{SR}}
\newcommand{\IC}{\mathrm{IC}}
\newcommand{\ICIR}{\mathrm{ICIR}}
\newcommand{\sgn}{\operatorname{sign}}
\newcommand{\rank}{\operatorname{rank}}
\newcommand{\zscore}{\operatorname{z}}

\title{\textbf{Alpha-GPT: Internal Research Note}}
\date{\today}

\begin{document}
\maketitle

\begin{abstract}
Finding a portfolio of alpha expressions that ranks future cross-sectional returns is a search
over an intractable, doubly-exponential space of programs against a black-box objective known
only by backtesting. We replace brute-force enumeration with \emph{economic-reasoning-guided
search}: a multi-agent LLM debate proposes economically motivated expressions (large, informed
jumps through the space), genetic programming refines them locally, and the debate's cross-review
plus a structural verifier prune weak candidates cheaply, so the expensive objective is queried
only on promising ones. The method specializes the \emph{system}, not the model -- a general LLM
already holds the asset-pricing reasoning we need, so we leave its weights untouched and wrap it
in a task-specific harness whose only levers are a decomposition (reason about the economic
mechanism, then commit to a formula) and test-time compute (the debate). We deliberately target
one modular stage of the alpha pipeline -- the generator -- holding data, portfolio construction,
and backtesting fixed: this attributes any measured edge to the alphas themselves and makes the
generator a drop-in module for existing quant infrastructure. Whether this latent economic skill
beats data-mining is the empirical question our random, single-shot, and genetic-programming
baselines are built to test.
\end{abstract}

% ===========================================================================
\section{Foundations: the objects we manipulate}
\label{sec:foundations}
% ===========================================================================

% ---------------------------------------------------------------------------
\subsection{Data: panels, universe, splits}
\label{sec:data}
% ---------------------------------------------------------------------------

The market is represented as a library of \emph{panels}. A panel for data field $k$
is a matrix
\[
X^{(k)} \in \R^{T \times N}, \qquad X^{(k)}_{t,i} = \text{value of field } k \text{ for stock } i \text{ on day } t,
\]
indexed by trading date $t$ (rows) and security identifier (PERMNO) $i$ (columns).
All fields share one date-by-stock grid, so they can be combined cell-wise. Missing
cells are \texttt{NaN} and every operator is NaN-safe.

The investable set on day $t$ is the point-in-time universe $\Uni_t \subseteq \{1,\dots,N\}$
(roughly 1{,}500 names per year, re-selected annually by trailing market
capitalization; survivorship-free, so names enter and exit at delisting). A cell is
valid only if $i \in \Uni_t$ and the stock actually traded that day.

\paragraph{The prediction target.} For each stock and day we define the
\emph{forward return}
\[
y_{t,i} \;=\; r_{t+1,i} \;=\; \frac{P_{t+1,i}}{P_{t,i}} - 1,
\]
the next-day total return (delisting returns embedded). Forward returns are masked to
the universe and the last in-sample observation is blanked so no label leaks across a
split boundary.

\paragraph{Splits.} The calendar is partitioned once:
\[
\underbrace{2010\text{--}2017}_{\text{train}} \;\big|\;
\underbrace{2018\text{--}2020}_{\text{validation}} \;\big|\;
\underbrace{2021\text{--}2025}_{\text{test}}.
\]
\emph{Selection of everything} (which expressions, which combiner, which weights, which
hyperparameters) happens on train/validation. The test split is touched exactly once,
by the final portfolio backtest. This constraint reappears in Section~\ref{sec:opt} as
an explicit feasibility condition on the optimizer.

\paragraph{Data sources and what we can add.} Table~\ref{tab:data} lists the live
sources and natural extensions. The key axes are \emph{source}, \emph{type},
\emph{native frequency}, and the orthogonal information each adds.

\begin{table}[h]
\centering
\small
\begin{tabular}{@{}p{2.6cm}p{3.2cm}p{1.9cm}p{5.3cm}@{}}
\toprule
Source & Type & Frequency & Notes / fields \\
\midrule
CRSP (CIZ) & price / volume & daily & OHLC, returns (delisting-adjusted), volume, shares, market cap, bid/ask. The price--volume base. \\
Compustat + WRDS ratios & fundamentals & quarterly & $\sim$17 characteristics (profitability, leverage, B/M, E/P, accruals, growth), point-in-time via earnings-announcement date \texttt{rdq}, forward-filled to daily. \\
\midrule
\multicolumn{4}{@{}l}{\emph{Natural extensions (orthogonal information families):}} \\
OSAP (Chen--Zimmermann) & published anomaly signals & monthly & $\sim$200 firm-level signals; ingestion is implemented but not currently loaded into the live universe. \\
\bottomrule
\end{tabular}
\caption{Data sources, types, native frequencies, and extensions. Lower-frequency
sources are forward-filled onto the daily grid with a point-in-time release lag, so
mixing frequencies never introduces look-ahead.}
\label{tab:data}
\end{table}

% ---------------------------------------------------------------------------
\subsection{An alpha expression}
\label{sec:alpha-expr}
% ---------------------------------------------------------------------------

\begin{definition}[Alpha expression]
An \emph{alpha expression} $e$ is a well-typed symbolic tree whose leaves are
\emph{terminals} (data panels, drawn from a finite menu $\Term$) and whose internal
nodes are \emph{operators} (drawn from a finite set $\Op$). Evaluating the tree
produces a \emph{signal panel}
\[
\alpha \;=\; e(X) \in \R^{T\times N}, \qquad
\alpha_t \;=\; \big(\alpha_{t,i}\big)_{i\in\Uni_t} \in \R^{N_t}
\]
where $\alpha_t$ is the day-$t$ cross-section of scores and $N_t = |\Uni_t|$.
\end{definition}

The operator set is small and fixed (the ``alpha language''):
\begin{itemize}[leftmargin=1.4em,itemsep=1pt]
  \item \textbf{Time-series} (per stock, along time; curried with fixed windows
  $w\in\{5,10,20,60\}$): \texttt{ts\_mean}, \texttt{ts\_std}, \texttt{ts\_delta},
  \texttt{ts\_rank}, \texttt{ts\_min}, \texttt{ts\_max}, \texttt{ts\_returns}, and the
  binary \texttt{ts\_corr}.
  \item \textbf{Cross-sectional} (per day, across stocks): \texttt{cs\_rank},
  \texttt{cs\_zscore}.
  \item \textbf{Element-wise}: \texttt{add}, \texttt{sub}, \texttt{mul},
  \texttt{safe\_div}, \texttt{log\_abs}, \texttt{abs\_val}, \texttt{sign},
  \texttt{neg}.
\end{itemize}
A concrete example, ``recent return scaled by volatility, ranked cross-sectionally,''
is
\[
e \;=\; \texttt{cs\_rank}\big(\texttt{safe\_div}(\texttt{ts\_returns\_20}(\texttt{close}),\ \texttt{ts\_std\_20}(\texttt{returns}))\big).
\]

\paragraph{Why the space is intractable.} Let $b$ be the typical operator arity and
$D$ the maximum tree depth. The number of admissible expressions grows roughly as
\begin{equation}
|\Expr| \;\sim\; \big(|\Term| + |\Op|\big)^{\,\Theta(b^{D})},
\label{eq:explosion}
\end{equation}
doubly exponential in depth. The search is over \emph{programs} (terminals, operators,
and nesting), not over feature pairs, so direct enumeration is infeasible and a guided
search is required.

\begin{remark}[Formula vs.\ code representation]
The tree above fixes one \emph{level} of representation -- a composition of the bounded operator
set $\Op$, i.e.\ a symbolic formula. A deeper alternative, central to concurrent work (CogAlpha;
Liu et al., 2026), lets the generator emit general \emph{code} (intermediate variables, regime
logic, transforms outside $\Op$), reaching the strictly larger
$\Expr_{\text{code}}\supset\Expr_{\text{formula}}$ and matching what LLMs do well. We default to
the DSL because it is fast, parseable, and open to the free genetic search of
Section~\ref{sec:numerics}, whereas code needs sandboxing, timeouts, and a look-ahead test,
evaluates more slowly, and forfeits free GP -- pushing toward far smaller candidate counts. Since
only this backend changes, ``generate code'' also drops in as a third arm for an all-else-equal
comparison the literature lacks.
\end{remark}

% ---------------------------------------------------------------------------
\subsection{From a signal to cross-sectional bets}
\label{sec:sorting}
% ---------------------------------------------------------------------------

A signal is used \emph{cross-sectionally}: each day we rank the stocks by $\alpha_t$
and bet that high-score names outperform low-score names. Formally, on day $t$ assign
each valid stock to one of $Q=5$ quantiles by the rank of its score,
\[
q_{t,i} \;=\; \Big\lceil Q \cdot \tfrac{\rank(\alpha_{t,i})}{N_t} \Big\rceil \in \{1,\dots,Q\},
\]
(ranking with tie-breaking so the bins never collapse). The long--short quintile
portfolio is long the top bucket and short the bottom bucket, equal-weighted, and its
realized next-day return is
\begin{equation}
R^{p}_{t}
\;=\;
\underbrace{\frac{1}{|L_t|}\sum_{i\in L_t} y_{t,i}}_{\text{long top quintile}}
\;-\;
\underbrace{\frac{1}{|S_t|}\sum_{i\in S_t} y_{t,i}}_{\text{short bottom quintile}},
\qquad
\begin{aligned}
L_t &= \{i: q_{t,i}=Q\},\\
S_t &= \{i: q_{t,i}=1\}.
\end{aligned}
\label{eq:lsret}
\end{equation}
Because positions are determined by \emph{rank}, only the cross-sectional ordering of
$\alpha_t$ matters, not its scale; any monotone transform of a signal is the same bet.

% ---------------------------------------------------------------------------
\subsection{Measuring one signal: the information coefficient}
\label{sec:ic}
% ---------------------------------------------------------------------------

The standard, scale-free measure of how well a signal ranks future returns is the
\emph{information coefficient}: the daily Spearman rank correlation between the score
and the forward return.

\begin{definition}[IC, ICIR]
For expression $e$ with signal $\alpha=e(X)$,
\begin{equation}
\IC_t(e) \;=\; \rho_{\text{Spearman}}\!\big(\alpha_t,\; y_t\big)
\;=\; \rho_{\text{Pearson}}\!\big(\rank(\alpha_t),\, \rank(y_t)\big),
\qquad
\IC(e) \;=\; \frac{1}{T}\sum_{t=1}^{T}\IC_t(e),
\label{eq:ic}
\end{equation}
and the information ratio of the signal is
\begin{equation}
\ICIR(e) \;=\; \frac{\overline{\IC_t}}{\,\mathrm{std}_t(\IC_t)\,},
\qquad
t\text{-stat} \;=\; \ICIR(e)\,\sqrt{T}.
\label{eq:icir}
\end{equation}
\end{definition}

$\IC$ measures average directional accuracy (a daily $\IC$ persistently above
$\approx 0.02$ is economically meaningful); $\ICIR$ measures its stability /
signal-to-noise. We also track \emph{turnover}, the mean day-over-day change in the
cross-sectional percentile rank,
\begin{equation}
\tau(e) \;=\; \frac{1}{T}\sum_{t} \ \frac{1}{N_t}\sum_i \big| u_{t,i} - u_{t-1,i}\big|,
\qquad u_{t,i} = \tfrac{\rank(\alpha_{t,i})}{N_t},
\label{eq:turnover}
\end{equation}
a proxy for rebalancing cost.

\begin{remark}[Why IC and not raw PnL for ranking candidates]
$\IC$ in \eqref{eq:ic} is a rank statistic, so it is robust to outliers and immune to
the monotone-transform ambiguity noted in Section~\ref{sec:sorting}. It is also cheap
to compute in bulk, which matters because the search evaluates enormous numbers of
candidates (Section~\ref{sec:numerics}).
\end{remark}

% ---------------------------------------------------------------------------
\subsection{Combining many signals into a portfolio}
\label{sec:combine}
% ---------------------------------------------------------------------------

A single alpha is weak; the edge comes from combining many weak,
weakly-correlated signals into a portfolio (this is the \emph{breadth} that
Section~\ref{sec:opt} shows matters most). Given a selected set $S=\{e_1,\dots,e_M\}$,
each signal is first \emph{oriented} and \emph{standardized}:
\begin{equation}
s^{(j)}_t \;=\; \sgn\!\big(\IC^{\text{val}}(e_j)\big)\cdot \zscore\!\big(\alpha^{(j)}_t\big),
\qquad
\zscore(\alpha_t)_i = \frac{\alpha_{t,i}-\mathrm{mean}(\alpha_t)}{\mathrm{std}(\alpha_t)},
\label{eq:orient}
\end{equation}
so every signal points the ``profitable'' way (sign chosen on validation, never test) and
lives on a comparable scale. The composite is then a weighted blend
$g_t = \sum_j w_j\, s^{(j)}_t$. Because portfolio construction is not the research question
(Section~\ref{sec:opt} expands on this), we use the simplest robust weights: the default is
\textbf{equal-weight} ($w_j = 1/M$), with \textbf{IC-weighting}
($w_j \propto |\IC^{\text{val}}(e_j)|$) as a variant. Both use only validation-derived signs
and weights, so no test information enters construction.

% ---------------------------------------------------------------------------
\subsection{The objective: out-of-sample Sharpe}
\label{sec:sharpe}
% ---------------------------------------------------------------------------

Running the composite $g$ through the long--short map \eqref{eq:lsret} yields a daily
portfolio return series $R^p_t$. The headline objective is its annualized Sharpe ratio
on the test split:
\begin{equation}
\boxed{\ \Sharpe \;=\; \sqrt{252}\;\frac{\mathrm{mean}_t\, R^p_t}{\mathrm{std}_t\, R^p_t}\ }
\label{eq:sharpe}
\end{equation}
We also report annualized return, maximum drawdown
$\mathrm{MDD}=\min_t\big(C_t/\max_{s\le t}C_s - 1\big)$ with $C_t=\prod_{s\le t}(1+R^p_s)$,
and the portfolio's own $\IC$/$\ICIR$/$\tau$. Equation~\eqref{eq:sharpe} is the scalar
that Part~I maximizes.

% ===========================================================================
\section{Part I: the optimization problem}
\label{sec:opt}
% ===========================================================================

The generator must choose a set $S\subseteq\Expr$ of expressions that, combined into a
long--short portfolio, maximizes out-of-sample Sharpe. Everything downstream is fixed by scope:
the combiner is not fitted, only picked from a small menu of robust weights on validation, so
there is no inner portfolio problem and $S$ is the sole decision variable:
\begin{equation}
\boxed{\
\max_{\,S \subseteq \Expr}\quad
\E\!\big[\Sharpe_{\text{test}}\!\big(\text{portfolio}(S)\big)\big]
\quad\text{s.t.}\quad
\begin{cases}
\text{every } e\in S \text{ is parseable and passes the verifier},\\
S \text{ and the combiner are chosen on train/validation only.}
\end{cases}
\ }
\label{eq:outer}
\end{equation}
The expectation matters because the optimizer never sees $\Sharpe_{\text{test}}$; it maximizes a
validation surrogate, and test is read exactly once. That validation--test gap is the central
statistical risk, reduced by the consistency filter (Section~\ref{sec:numerics}) and the debate's
cross-review. The problem is hard for two reasons -- $\Expr$ is doubly exponential
\eqref{eq:explosion} and each evaluation is an expensive backtest -- so we can neither enumerate
nor differentiate, only \emph{search} (Part~II).

\paragraph{What the objective rewards: skill and breadth.}
For a dollar-neutral long--short book the portfolio return \emph{is} the active return, so its
Sharpe equals the \emph{information ratio}, and Grinold's fundamental law factors that same number
into the two levers a proposer controls:
\begin{equation}
\boxed{\ \Sharpe \;=\; \mathrm{IR} \;\approx\; \underbrace{\IC}_{\text{skill}}\;\cdot\;\sqrt{\,\underbrace{\mathrm{breadth}}_{\#\ \text{independent bets}}\,}\ }
\label{eq:funlaw}
\end{equation}
Skill enters linearly but breadth only under a square root, so many diverse, modest-$\IC$ alphas
beat a few strong ones. The search therefore maximizes the \emph{number of independent} skilled
alphas rather than individual quality, and the combiner keeps $S$ large with no top-$k$ cap --
provided the added alphas are genuinely independent, not restatements of the same bet.
``Skill $\times$ breadth'' is also the axis along which the proposal mechanisms of Part~II differ
(Section~\ref{sec:whydebate}).

\begin{remark}[Combiner: equal weight, not fitted tangency]
The Sharpe-maximizing weights for a fixed $S$ are the tangency portfolio
$w^\star \propto \Sigma^{-1}\mu$. With many correlated, noisy alphas the sample $\Sigma^{-1}$ is
ill-conditioned and $w^\star$ overfits; equal weight is its maximal-shrinkage limit, beats sample
tangency out of sample, and keeps any measured edge attributable to the alphas rather than to
combiner tuning.
\end{remark}

% ===========================================================================
\section{Part II: how the search works}
\label{sec:search}
% ===========================================================================

Since \eqref{eq:outer} can be neither enumerated nor differentiated, we search. The search is
a generate-and-test loop built to spend cheap effort widely and the expensive true objective
narrowly: two complementary moves generate candidates, and a cascade of increasingly expensive
checks filters them, so the real backtest runs only on the few that survive.

% ---------------------------------------------------------------------------
\subsection{Propose globally, refine locally, prune cheaply}
\label{sec:approx}
% ---------------------------------------------------------------------------

\begin{enumerate}[leftmargin=1.6em,itemsep=3pt]
  \item \textbf{Propose (global moves).} The multi-agent debate takes an economically reasoned
  hypothesis and writes \emph{complete} expressions for it. Because it draws on the model's
  pretrained economic knowledge, it lands directly in plausible regions of $\Expr$ instead of
  wandering in from random starts -- large, informed jumps.
  \item \textbf{Refine (local moves).} Each proposal seeds genetic programming, which mutates
  and recombines it (subtree crossover, point mutation) to hill-climb train $\IC$ in a small
  neighborhood. One informed seed thus spawns many nearby variants at no extra model cost --
  small, cheap steps.
  \item \textbf{Prune (cheap first, expensive last).} Running the true backtest is the costly
  step, so cheaper checks come first, in order: the debate's cross-review drops economically
  weak ideas, the structural verifier discards anything that fails to parse, evaluate, or vary,
  and a fast approximate $\IC$ ranks what remains. Only the top survivors reach the full
  backtest.
\end{enumerate}

The division of labor is the point: the LLM makes the large, economically informed jumps a
local search never could, and GP does the cheap refinement around them that the LLM would waste
calls doing by hand. The proposer is \emph{fixed} -- it does not learn from which past alphas
backtested well.

% ---------------------------------------------------------------------------
\subsection{Why multi-agent debate}
\label{sec:whydebate}
% ---------------------------------------------------------------------------

\paragraph{Diversity, then critique.} Search quality is set by the \emph{best} candidate found
under a fixed proposal budget. One model sampled many times tends to return the same few ideas
reworded, so most of the budget re-draws a single mode; three persona-conditioned proposers
(Momentum, MeanReversion, Fundamental) instead pull toward \emph{different} economic mechanisms,
so the same budget covers more of $\Expr$ and lifts the expected best. Debate then adds what a
plain sampler cannot: each agent critiques the others (cross-review), killing weak proposals
before any of them costs a backtest. Debate thus buys both \emph{exploration} (diverse
proposals) and \emph{verification} (critique).

\paragraph{Why debate and not the other inference-time methods.}
\begin{itemize}[leftmargin=1.4em,itemsep=2pt]
  \item \textbf{Self-consistency / majority vote} collapses many reasoning chains onto one
  consensus answer; we want the opposite -- a \emph{diverse set} of distinct alphas.
  \item \textbf{Best-of-$N$} samples one model and keeps the top-scored; with no persona
  diversity and no critique, its proposals cluster and it explores narrowly.
  \item \textbf{MCTS / tree-of-thought} needs a cheap, accurate value estimate for partial
  programs to steer the tree; here the only ground truth is the expensive backtest, so there is
  nothing cheap to steer by.
\end{itemize}

\paragraph{The hypothesis: debate is the only mechanism with both skill and breadth.} The
fundamental law \eqref{eq:funlaw} says a good proposer needs \emph{skill} (high per-alpha $\IC$)
and \emph{breadth} (many independent alphas), and the skill that counts is \emph{out-of-sample}.
This is exactly where an economic prior earns its keep: brute-force enumeration and genetic
programming can push up \emph{in-sample} $\IC$, but with no economic reasoning to anchor them
their best fits are usually overfit and fade out of sample. Table~\ref{tab:skillbreadth} lays out
where we \emph{expect} each mechanism to land: the priorless methods fall short on genuine skill,
breadth, or both, whereas debate is built for both -- genuine skill from economic priors plus
cross-review, breadth from cheap, diverse proposals. These are predictions; the ablation grid of
Section~\ref{sec:baselines} measures whether they hold.

\begin{table}[h]
\centering
\small
\begin{tabular}{@{}p{2.9cm}p{1.4cm}p{1.7cm}p{7.4cm}@{}}
\toprule
Mechanism & Skill & Breadth & Why we expect this \\
\midrule
Random trees        & none     & high     & unlimited cheap trees, but no economic prior -- the signals are meaningless \\
Brute-force enum.\  & low      & very low & no economic prior, so its high in-sample $\IC$ is mostly overfit; also reaches only a shallow sliver of $\Expr$ \\
Genetic programming & low      & low      & hill-climbs in-sample $\IC$ with no economic grounding, so its winners tend to overfit; explores only locally \\
Single-shot LLM     & moderate & moderate & an economic prior gives a genuine but unrefined edge; one call clusters on a theme and gets no critique \\
Multi-agent debate  & high     & high     & economic priors plus cross-review favor genuinely predictive alphas; diverse personas widen coverage \\
\bottomrule
\end{tabular}
\caption{\emph{Hypothesized} skill and breadth by proposal mechanism -- predictions, not
measurements. \emph{Skill} is genuine \emph{out-of-sample} per-alpha $\IC$: the priorless
mechanical methods can fit the training data, but with no economic reasoning to lean on that fit
tends to overfit rather than generalize. \emph{Breadth} is the number of independent alphas the
mechanism can supply. Only debate is designed for both. Section~\ref{sec:baselines} tests these
predictions.}
\label{tab:skillbreadth}
\end{table}

% ---------------------------------------------------------------------------
\subsection{Baselines and the ablation grid}
\label{sec:baselines}
% ---------------------------------------------------------------------------

To test these predictions we run the debate against baselines that strip out the economic prior
and the multi-agent structure, each through an \emph{identical} downstream pipeline (verify,
evaluate, filter, combine, backtest; Part~III), so any difference traces to the proposal mechanism
alone. Crossing the mechanism with a GP on/off toggle gives the grid in Table~\ref{tab:ablation}:
reading down a column isolates the economic prior (random vs.\ single-shot LLM vs.\ debate),
reading across a row isolates local GP amplification (off vs.\ on).

\begin{table}[h]
\centering
\small
\begin{tabular}{@{}llll@{}}
\toprule
Proposal mechanism & GP off & GP on & LLM cost / round \\
\midrule
Random trees (noise floor)      & \texttt{random}        & \texttt{random-gp}        & free \\
Single-shot LLM (one call)      & \texttt{single-agent}  & \texttt{single-agent-gp}  & $\sim$1 call \\
Multi-agent debate (the method) & \texttt{debate-only}   & \texttt{full}             & $\sim$10--12 calls \\
\bottomrule
\end{tabular}
\caption{Ablation grid: proposal mechanism $\times$ GP toggle (cell labels are the implemented
run modes). The downstream verify / evaluate / filter / combine / backtest path is identical
across all six cells, so differences come only from how candidates are proposed and amplified.}
\label{tab:ablation}
\end{table}

The down-column contrast -- does an economic prior, and \emph{multi-agent} structure, beat the
mechanical baselines? -- is the project's central hypothesis; the sharpest single test is
\texttt{single-agent-gp} vs.\ \texttt{debate-only}, which at matched downstream effort separates
structured multi-agent reasoning from merely spending more compute on one proposer.

\begin{remark}[Debate $+$ GP]
Adding GP on top of the debate pushes breadth further -- each economically-reasoned seed spawns
many verified neighbors -- but GP's crossover and mutation are economically blind, so the
variants can dilute the skill that made the seeds good. Debate $+$ GP thus trades a little
per-alpha skill for breadth; whether that is net-positive is what the GP-on column of
Table~\ref{tab:ablation} answers.
\end{remark}

% ===========================================================================
\section{Part III: numerics and algorithms}
\label{sec:numerics}
% ===========================================================================

The concepts of Part~II become three algorithms operating on two objects -- the \emph{context}
$\mathcal{H}$ (the hypothesis under exploration plus the resolved terminal menu $\Term$) and the
\emph{book} $\mathcal{P}$ (every expression found so far with its cached train/validation
metrics) -- and four components: three \textbf{debate agents} (Momentum, MeanReversion,
Fundamental), a \textbf{moderator} that synthesizes their proposals, a \textbf{GP operator} that
mutates and recombines seeds, and a \textbf{verifier} that structurally filters candidates
(parse, depth, eval, coverage, non-constant) to keep the book clean at scale.

% ---------------------------------------------------------------------------
\subsection{The multi-agent debate (the proposer)}
% ---------------------------------------------------------------------------

Each of the two stages is a draft / cross-review / revise loop over the three agents,
closed by a moderator. Stage 1 produces hypotheses (no formulas yet); Stage 2 turns
them into parseable expressions.

\begin{algorithm}[h]
\caption{Two-stage debate: idea $\to$ formula (one full proposal round)}
\begin{algorithmic}[1]
\Require trading idea / context $\mathcal{H}$, terminal menu $\Term$, agents $\{1,2,3\}$, moderator $\Phi$
\Statex \textbf{Stage 1: idea debate}
\State each agent $i$ \textbf{drafts} a structured hypothesis (mechanism, direction, proxies, filters, normalization)
\State each agent \textbf{cross-reviews} the others on 8 dimensions (mechanism, clarity, payoff, direction, subfactor, filter, normalization, implementability), decision $\in$ \{accept, accept-with-revision, reject\}
\State each agent \textbf{revises} its hypothesis given the reviews
\State $\mathcal{H}^\star \gets \Phi(\text{revised hypotheses})$ \Comment{moderator synthesizes $\approx 3$ hypotheses}
\Statex \textbf{Stage 2: formula debate}
\State each agent \textbf{drafts} 1--2 expressions per hypothesis over $\Term$
\State each agent \textbf{cross-reviews} on \{faithfulness, implementability, robustness, novelty, simplicity\}
\State each agent \textbf{revises}; \textbf{parse-gate}: drop any expression that does not compile against $(\Term,\Op)$
\State \Return seed pack $\gets \Phi(\text{parseable expressions})$ \Comment{moderator selects $\approx 8$ diverse seeds}
\end{algorithmic}
\end{algorithm}

Cost is roughly 10--12 model calls per round (3 agents $\times$ 3 phases $\times$ 2
stages plus 2 moderator calls). This is the expensive, high-quality proposal path; the
factory can also use a single autonomous generator call when scale matters more than
per-idea quality.

% ---------------------------------------------------------------------------
\subsection{The outer search loop (factory + GP amplification)}
% ---------------------------------------------------------------------------

\begin{algorithm}[h]
\caption{Outer search: propose $\to$ amplify $\to$ verify $\to$ evaluate $\to$ store}
\begin{algorithmic}[1]
\Require number of rounds $N$, train/val splits, primitive set $\mathrm{pset}$
\State $\mathcal{P}\gets\emptyset$ \Comment{the book / database, deduplicated on expression hash}
\For{$n = 1,\dots,N$}
  \State seeds $\gets \textsc{Propose}(\mathcal{H}_n)$ \Comment{Algorithm 1, or single-agent / random}
  \State candidates $\gets$ seeds $\cup\ \textsc{GP-amplify}(\text{seeds})$ \Comment{local search around each seed}
  \For{each $e$ in candidates}
    \If{\textbf{not} \textsc{Verify}$(e)$} \textbf{continue} \Comment{parse, depth, eval, coverage, non-constant}
    \EndIf
    \State compute $\IC^{\text{train}}(e), \IC^{\text{val}}(e), \ICIR, \Sharpe^{\text{val}}, \tau$ \Comment{fast approximate metrics}
    \State insert $e$ into $\mathcal{P}$
  \EndFor
\EndFor
\State \Return $\mathcal{P}$
\end{algorithmic}
\end{algorithm}

\textsc{GP-amplify} runs DEAP genetic programming seeded with the parsed proposals (subtree
crossover, point mutation, depth limit 6; fitness $=$ mean train $\IC$ on a date-subsample for
speed), so each informed seed spawns many verified neighbors with no extra model calls.

% ---------------------------------------------------------------------------
\subsection{Selection and the out-of-sample backtest}
% ---------------------------------------------------------------------------

\begin{algorithm}[h]
\caption{From the book to one out-of-sample number}
\begin{algorithmic}[1]
\Require book $\mathcal{P}$, splits (val, test), combiners $\mathcal{C}$
\State $S \gets \{e\in\mathcal{P} : \IC^{\text{train}}(e)\cdot \IC^{\text{val}}(e) > 0 \ \text{and}\ |\IC^{\text{val}}(e)| \ge \epsilon\}$ \Comment{consistency filter, no top-$k$ cap}
\State orient + z-score each $e\in S$ on test via \eqref{eq:orient} \Comment{signs from validation}
\For{each combiner $c\in\mathcal{C}$}
  \State $g \gets$ combine$(S, c)$; \quad $R^p \gets$ long--short quintile of $g$ on test \Comment{\eqref{eq:lsret}}
  \State record $\Sharpe, \IC, \ICIR, \tau, \mathrm{MDD}$ \Comment{\eqref{eq:sharpe}}
\EndFor
\State choose the headline combiner by \textbf{validation} Sharpe \Comment{not by max test Sharpe}
\State \Return the test metrics of the chosen combiner
\end{algorithmic}
\end{algorithm}

\paragraph{Consistency filter, not top-$k$.} Selecting the highest-validation-$\IC$ alphas tends
to pick the most overfit ones; instead we require the sign to agree across train and validation
($\IC^{\text{train}}\!\cdot\IC^{\text{val}}>0$), which drops the lucky ones while keeping the set
large. Even the combiner is chosen on validation, so no decision peeks at test.

% ---------------------------------------------------------------------------
\subsection{Evaluation speed and cost}
% ---------------------------------------------------------------------------

$\IC$ is always the exact per-day Spearman \eqref{eq:ic} -- one metric used both to rank
candidates during the search and to report the final portfolio. Sharpe comes in two forms: a
fast rank-weighted proxy for ranking candidates, and the exact quintile long--short backtest
\eqref{eq:lsret} for the headline number, run only on the filtered survivors. Everything
parallelizes; at roughly \$0.00008 per alpha, $10^3$ alphas cost about \$0.06 in $\approx$13
minutes and scale to $10^6$, so wall-clock, not money, binds -- which is why cheap GP
amplification per expensive proposal matters.

\end{document}
